% ============================================================
% Model: 量子二能级系统
% 覆盖：中微子振荡、自旋在磁场中、拉比振荡、微扰论
% ============================================================

\section{量子二能级系统}
\label{sec:two-level-system}

\begin{tipbox}[关键词标签]
二能级系统 · Bloch向量 · Bloch球 · 微扰论 · 拉比振荡 · 中微子振荡 · 混合角 · 含时演化 · 能量修正
\end{tipbox}

% --------------------------------------------------------
% 1. 模型定义框
% --------------------------------------------------------
\begin{modelbox}[模型：$H_0$ 本征基下的二能级系统]
\textbf{物理情境}：许多量子系统在某个能量尺度下可近似为只有两个相关能级参与的过程。典型的二能级系统包括：中微子振荡（电子中微子 vs 缪子中微子）、自旋-$1/2$ 在磁场中（自旋向上 vs 向下）、氨分子反转（氮在氢平面之上 vs 之下）、量子比特（$|0\rangle$ vs $|1\rangle$）。

\textbf{通用哈密顿量}：
\[
H=H_0+V=\begin{pmatrix}E_1&V_{12}\\V_{21}&E_2\end{pmatrix},
\quad V_{12}=V_{21}^*.
\]

\textbf{微扰论框架}（$|V_{12}|\ll|E_1-E_2|$）：
\begin{itemize}[nosep]
    \item 未微扰本征态：$|1\rangle,|2\rangle$，能量 $E_1^{(0)},E_2^{(0)}$
    \item 一阶能量修正：$E_n^{(1)}=\langle n|V|n\rangle=V_{nn}$
    \item 二阶能量修正：$E_n^{(2)}=\sum_{m\neq n}\frac{|\langle m|V|n\rangle|^2}{E_n^{(0)}-E_m^{(0)}}$
\end{itemize}

\textbf{精确对角化}（$2\times2$ 矩阵）：
\begin{itemize}[nosep]
    \item 定义平均能量 $\bar{E}=(E_1+E_2)/2$，能级分裂 $\Delta=E_2-E_1$
    \item 定义有效分裂 $\Omega=\sqrt{\Delta^2+4|V_{12}|^2}$
    \item 本征值：$E_{\pm}=\bar{E}\pm\frac{\Omega}{2}$
    \item 混合角 $\theta$：$\tan2\theta=\frac{2|V_{12}|}{\Delta}$
\end{itemize}
\end{modelbox}

% --------------------------------------------------------
% 2. 关键公式框
% --------------------------------------------------------
\begin{formulabox}[核心公式]
\begin{enumerate}[label=\textbf{(\arabic*)}]
    \item \textbf{微扰论能量修正（非简并）}：
    \[
    E_n\approx E_n^{(0)}+\langle n|V|n\rangle+\sum_{m\neq n}\frac{|\langle m|V|n\rangle|^2}{E_n^{(0)}-E_m^{(0)}}+\cdots
    \]
    \texteq{适用条件：$|\langle m|V|n\rangle|\ll|E_n^{(0)}-E_m^{(0)}|$}

    \item \textbf{二能级精确本征值}：
    \[
    E_{\pm}=\frac{E_1+E_2}{2}\pm\frac{1}{2}\sqrt{(E_2-E_1)^2+4|V_{12}|^2}
    \]
    \texteq{$E_+-E_-=\Omega=\sqrt{\Delta^2+4|V_{12}|^2}$}

    \item \textbf{混合角与精确本征态}：
    \[
    |+\rangle=\cos\theta\,|1\rangle+\sin\theta\,|2\rangle,\quad
    |-\rangle=-\sin\theta\,|1\rangle+\cos\theta\,|2\rangle,
    \]
    \[
    \tan2\theta=\frac{2|V_{12}|}{E_2-E_1},\quad 0\leq\theta\leq\frac{\pi}{2}.
    \]

    \item \textbf{含时演化与 survival probability}：
    初态 $|\psi(0)\rangle=|1\rangle$，则
    \[
    P_{1\to1}(t)=|\langle1|\psi(t)\rangle|^2=1-\sin^2(2\theta)\,\sin^2\Bigl(\frac{\Omega t}{2\hbar}\Bigr).
    \]
    \texteq{振荡频率由能级分裂 $\Omega$ 决定，振幅由混合角决定}

    \item \textbf{拉比振荡}（自旋在磁场中，$H=-B_z\sigma_z-\mu\sigma_x$）：
    \[
    P_{\uparrow}(t)=\frac{\mu^2}{\mu^2+B_z^2}\sin^2\Bigl(\frac{\sqrt{\mu^2+B_z^2}}{\hbar}t\Bigr)
    \]
    \texteq{初态 $|\downarrow\rangle$，$B_z$ 为纵向场，$\mu$ 为横向扰动强度}

    \item \textbf{中微子振荡概率}：
    \[
    P_{e\to e}(L)=1-\sin^2(2\theta)\,\sin^2\Bigl(\frac{\Delta m^2\,L}{4E\hbar c}\Bigr)
    \]
    \texteq{$\Delta m^2=m_2^2-m_1^2$，$\theta$ 为 Pontecorvo-Maki-Nakagawa-Sakata 混合角}

    \item \textbf{振荡周期与振荡长度}：
    \[
    T_{\text{osc}}=\frac{4\pi\hbar}{\Omega},\qquad
    L_{\text{osc}}=\frac{4\pi E\hbar c}{\Delta m^2}.
    \]
\end{enumerate}
\end{formulabox}

% --------------------------------------------------------
% 3. 训练点框
% --------------------------------------------------------
\begin{drillbox}[训练点与考法变体]
\trainingpoint{1} \textbf{微扰论适用条件判断}：给定 $H_0$ 和 $V$，判断是否可用微扰论。关键：$|V_{nm}|\ll|E_n^{(0)}-E_m^{(0)}|$。若近简并，微扰论失效，需严格对角化。

\trainingpoint{2} \textbf{二阶能量修正计算}：给定 $V$ 和一组未微扰态，写出二阶修正的完整求和表达式。注意分母符号：$E_n^{(0)}-E_m^{(0)}$，若 $m$ 在 $n$ 之上，分母为负。

\trainingpoint{3} \textbf{精确解 vs 微扰解对比}：对给定 $2\times2$ 矩阵，同时计算精确本征值和微扰论近似值，在 $|V_{12}|\ll\Delta$ 极限下展开精确解至 $O(V^2)$，验证与微扰论一致。

\trainingpoint{4} \textbf{拉比振荡参数提取}：给定 $P_{\uparrow}(t)$ 的实验数据（振荡幅度和频率），反推 $B_z$ 和 $\mu$。幅度 $A=\mu^2/(\mu^2+B_z^2)$，频率 $\omega=2\sqrt{\mu^2+B_z^2}/\hbar$。

\trainingpoint{5} \textbf{中微子振荡参数计算}：给定基线 $L$、能量 $E$ 和观测到的振荡概率，求 $\Delta m^2$ 和 $\sin^2(2\theta)$。注意概率公式的范围 $P\in[1-\sin^2(2\theta),1]$。

\trainingpoint{6} \textbf{共振条件}：当 $E_1=E_2$（简并或近简并），$\theta=\pi/4$，混合最大。此时 $P_{1\to1}(t)=\cos^2(\Omega t/2\hbar)$，粒子在两态之间完全来回振荡。

\trainingpoint{7} \textbf{Bloch 向量与态的对应}：给定 Bloch 向量 $\mathbf{n}$，反求密度矩阵、判断纯/混合态、计算各自旋分量的期望值。给定纯态 $|\psi\rangle$，求其在 Bloch 球上的坐标 $(\theta,\phi)$。

\trainingpoint{8} \textbf{Bloch 方程的几何解}：$d\mathbf{n}/dt=\boldsymbol{\Omega}\times\mathbf{n}$ 描述 Bloch 向量绕 $\boldsymbol{\Omega}$ 轴的匀速进动。初态在 $x$ 轴（$\mathbf{n}=(1,0,0)$），$\boldsymbol{\Omega}=(0,0,\omega)$，求 $\mathbf{n}(t)$。

\trainingpoint{9} \textbf{二能级系统的绝热演化}：若参数随时间缓慢变化（如磁场方向缓慢旋转），系统始终保持在瞬时本征态上（Berry 相位）。
\end{drillbox}

% --------------------------------------------------------
% 4. 解题技巧框
% --------------------------------------------------------
\begin{tipbox}[快速识别与解题技巧]
\begin{itemize}
    \item \examnote{识别标志}：题干出现``二能级''''微扰论''''拉比振荡''''中微子振荡''''自旋磁场''''混合角''，立即定位本模型。
    \item \examnote{$2\times2$ 矩阵速算}：
    \begin{center}
    \fbox{\parbox{0.9\linewidth}{\centering
    $H=\begin{pmatrix}a&b\\b^*&c\end{pmatrix}$，$a,c\in\mathbb{R}$\\[4pt]
    本征值：$\lambda_{\pm}=\frac{a+c}{2}\pm\sqrt{\bigl(\frac{a-c}{2}\bigr)^2+|b|^2}$\\[4pt]
    混合角：$\tan2\theta=\frac{2|b|}{c-a}$，$|+\rangle=\cos\theta|1\rangle+\sin\theta|2\rangle$
    }}
    \end{center}
    \item \examnote{微扰论 vs 精确解速查}：
    \begin{center}
    \begin{tabular}{lll}
    \toprule
    \textbf{量} & \textbf{微扰论} & \textbf{精确解} \\
    \midrule
    $E_+$ & $E_1+V_{11}+\frac{|V_{12}|^2}{E_1-E_2}$ & $\bar{E}+\frac{\Omega}{2}$ \\
    $E_-$ & $E_2+V_{22}+\frac{|V_{12}|^2}{E_2-E_1}$ & $\bar{E}-\frac{\Omega}{2}$ \\
    $\theta$ & $\frac{|V_{12}|}{E_2-E_1}$（小角） & $\frac{1}{2}\arctan\frac{2|V_{12}|}{\Delta}$ \\
    \bottomrule
    \end{tabular}
    \end{center}
    \item \examnote{Bloch 球速查}：
    \begin{center}
    \begin{tabular}{lll}
    \toprule
    \textbf{态} & \textbf{Bloch 向量} & \textbf{球面位置} \\
    \midrule
    $|1\rangle$ & $(0,0,1)$ & 北极 \\
    $|2\rangle$ & $(0,0,-1)$ & 南极 \\
    $|+\rangle=\frac{1}{\sqrt2}(|1\rangle+|2\rangle)$ & $(1,0,0)$ & $+x$ 轴 \\
    $|-\rangle=\frac{1}{\sqrt2}(|1\rangle-|2\rangle)$ & $(-1,0,0)$ & $-x$ 轴 \\
    $|+i\rangle=\frac{1}{\sqrt2}(|1\rangle+i|2\rangle)$ & $(0,1,0)$ & $+y$ 轴 \\
    完全混合 & $(0,0,0)$ & 球心 \\
    \bottomrule
    \end{tabular}
    \end{center}
    \item \examnote{概率守恒检查}：$P_{1\to1}(t)+P_{1\to2}(t)=1$。含时演化中总概率必须守恒，这是检验计算正确性的重要手段。
    \item \examnote{中微子振荡的极端相对论近似}：$E_k=\sqrt{p^2c^2+m_k^2c^4}\approx pc+\frac{m_k^2c^3}{2p}$，故 $E_2-E_1\approx\frac{\Delta m^2c^3}{2p}\approx\frac{\Delta m^2c^4}{2E}$。
    \item \examnote{拉比公式记忆}：把 $H=-B_z\sigma_z-\mu\sigma_x$ 写成 $H=-\mathbf{B}\cdot\boldsymbol{\sigma}$，其中 $\mathbf{B}=(\mu,0,B_z)$。有效磁场大小 $B_{\text{eff}}=\sqrt{\mu^2+B_z^2}$。拉比频率 $\omega_R=2B_{\text{eff}}/\hbar$。
\end{itemize}
\end{tipbox}

% --------------------------------------------------------
% 5. 易错点框
% --------------------------------------------------------
\begin{errorbox}{常见失误与陷阱}
\begin{itemize}
    \item \textbf{微扰论适用条件的误判}：微扰论要求 $|V_{12}|\ll|E_1-E_2|$。若两能级近简并，分母趋于零，二阶修正发散，必须严格对角化。常见错误是在近简并时仍用微扰论。
    \item \textbf{二阶修正分母符号}：$E_n^{(2)}=\sum_{m\neq n}\frac{|\langle m|V|n\rangle|^2}{E_n^{(0)}-E_m^{(0)}}$。若 $m$ 在 $n$ 之上，分母为负，能量修正为负（基态被压低）；若 $m$ 在 $n$ 之下，分母为正，能量修正为正。这与``所有能级相互排斥''的物理图像一致。
    \item \textbf{混合角的范围}：$\tan2\theta=\frac{2|V_{12}|}{E_2-E_1}$。若 $E_2<E_1$，则 $\tan2\theta<0$，$\theta\in(\pi/4,\pi/2]$。此时 $|+\rangle$ 主要是 $|1\rangle$ 成分，$|-\rangle$ 主要是 $|2\rangle$ 成分——能量高的态对应``高成分''的本征态。
    \item \textbf{拉比振荡的振幅}：$P_{\uparrow}(t)$ 的最大值是 $\frac{\mu^2}{\mu^2+B_z^2}$，不是 $1$（除非 $B_z=0$）。当 $B_z\gg\mu$ 时，振幅 $\approx(\mu/B_z)^2\ll1$，跃迁被抑制。
    \item \textbf{时间演化中的相位}：$|\psi(t)\rangle=c_1(t)e^{-iE_1t/\hbar}|1\rangle+c_2(t)e^{-iE_2t/\hbar}|2\rangle$。若用相互作用绘景，去掉自由演化相位后，$c_1(t)$ 和 $c_2(t)$ 的演化由 $V$ 决定。两种绘景的结果一致，但中间步骤不同，不要混用。
    \item \textbf{中微子振荡与 $L/E$}：振荡概率只依赖于比值 $L/E$（在固定 $\Delta m^2$ 下）。实验上通过改变基线 $L$ 或中微子能量 $E$ 来测量振荡。
\end{itemize}
\end{errorbox}

% --------------------------------------------------------
% 6. 物理图像与关联模型
% --------------------------------------------------------
\begin{insightbox}[物理图像与模型关联]
\textbf{物理图像}：

二能级系统就像一个跷跷板：
\begin{itemize}[nosep]
    \item 未微扰时（$V=0$），两个能级 $|1\rangle$ 和 $|2\rangle$ 各自独立，像跷跷板两端平衡着不同重量的人。
    \item 微扰 $V$ 像连接两端的弹簧。弹簧越弱，两端越独立；弹簧越强，两端耦合越紧。
    \item 当弹簧很弱（$|V_{12}|\ll\Delta$），两端几乎不动，只是轻微振动——这就是微扰论。
    \item 当弹簧很强或两端重量相等（简并），弹簧使两端强烈耦合，跷跷板剧烈振荡——这就是拉比振荡或中微子振荡。
    \item 混合角 $\theta$ 描述了``每个本征态中有多少成分是 $|1\rangle$、多少是 $|2\rangle$''。$\theta=\pi/4$ 时完全混合（各 $50\%$）。
\end{itemize}

\textbf{多种物理情境的统一描述}：
\begin{center}
\begin{tabular}{llll}
\toprule
\textbf{物理系统} & \textbf{``态 1''} & \textbf{``态 2''} & \textbf{微扰/耦合} \\
\midrule
中微子振荡 & $|\nu_e\rangle$ & $|\nu_\mu\rangle$ & 质量矩阵非对角元 \\
自旋磁场 & $|\uparrow\rangle$ & $|\downarrow\rangle$ & 横向磁场 $\mu\sigma_x$ \\
氨分子 & 氮在上 & 氮在下 & 隧穿势垒 \\
量子比特 & $|0\rangle$ & $|1\rangle$ & 微波驱动场 \\
NMR & 自旋平行 & 自旋反平行 & 射频脉冲 \\
\bottomrule
\end{tabular}
\end{center}

\textbf{与相似模型的对比}：
\begin{center}
\begin{tabular}{llll}
\toprule
\textbf{对比维度} & \textbf{微扰论} & \textbf{精确对角化} & \textbf{含时微扰} \\
\midrule
哈密顿量 & $H=H_0+V$（$V$ 小） & $H$（任意 $2\times2$） & $H(t)=H_0+V(t)$ \\
本征态 & 未微扰态 & 精确本征态 & 瞬时本征态 \\
能量 & 修正级数 & 精确值 & 无定态 \\
时间演化 & 相位因子 $e^{-iEt/\hbar}$ & 相位因子 & 跃迁概率（Fermi黄金律） \\
\bottomrule
\end{tabular}
\end{center}

\textbf{推广方向}：
\begin{itemize}[nosep]
    \item 三能级系统（如 $K^0$-$\bar{K}^0$ 系统）：引入 CP 破坏参数 $\epsilon$
    \item 含时驱动（Floquet 理论）：周期性驱动下的准能级
    \item 开放量子系统： Lindblad 主方程描述退相干对振荡的阻尼
    \item Berry 相位：绝热循环演化中积累的拓扑相位
    \item 量子 Zeno 效应：频繁测量抑制态的跃迁
\end{itemize}

\textbf{Bloch 球的物理图像}：

Bloch 球是二能级系统的``状态地图''：
\begin{itemize}[nosep]
    \item 球面上的每一点对应一个纯态，球内的点对应混合态。
    \item 北极是 $|1\rangle$，南极是 $|2\rangle$，赤道上的点是 $|1\rangle$ 和 $|2\rangle$ 的等概率叠加（相位决定具体位置）。
    \item 哈密顿量的非对角元 $V_{12}$ 产生绕 $x$ 或 $y$ 轴的转动；对角分裂 $\Delta$ 产生绕 $z$ 轴的转动。
    \item 拉比振荡就是 Bloch 向量在球面上画出一个圆（若 $|\boldsymbol{\Omega}|$ 恒定）。
    \item 退相干使 Bloch 向量从球面``缩向''球心——纯态变为混合态。
\end{itemize}
\end{insightbox}

% --------------------------------------------------------
% 7. 推导附录
% --------------------------------------------------------
\begin{derivationbox}[推导细节（自我检验用）]
\textbf{Step 1：二能级精确对角化}

哈密顿量：
\[
H=\begin{pmatrix}E_1&V\\V^*&E_2\end{pmatrix},\quad V\equiv V_{12}.
\]

本征方程 $\det(H-\lambda I)=0$：
\[
(E_1-\lambda)(E_2-\lambda)-|V|^2=0.
\]

解：
\[
\lambda_{\pm}=\frac{E_1+E_2}{2}\pm\frac{1}{2}\sqrt{(E_2-E_1)^2+4|V|^2}
=\bar{E}\pm\frac{\Omega}{2}.
\]

本征态：设 $|+\rangle=\cos\theta\,|1\rangle+\sin\theta\,|2\rangle$，代入 $H|+\rangle=E_+|+\rangle$：
\[
\begin{cases}
E_1\cos\theta+V\sin\theta=E_+\cos\theta,\\
V^*\cos\theta+E_2\sin\theta=E_+\sin\theta.
\end{cases}
\]

由第一式：
\[
V\sin\theta=(E_+-E_1)\cos\theta=\Bigl(\frac{\Omega-\Delta}{2}\Bigr)\cos\theta.
\]

利用 $\Omega^2=\Delta^2+4|V|^2$，得 $\Omega-\Delta=\frac{4|V|^2}{\Omega+\Delta}$，故：
\[
\tan\theta=\frac{\Omega-\Delta}{2V}=\frac{2|V|}{\Omega+\Delta}
\quad\Rightarrow\quad
\tan2\theta=\frac{2|V|}{\Delta}.
\]

\textbf{Step 2：含时演化与 survival probability}

初态 $|\psi(0)\rangle=|1\rangle$。用精确本征态展开：
\[
|1\rangle=\cos\theta\,|+\rangle-\sin\theta\,|-\rangle.
\]

时间演化：
\[
|\psi(t)\rangle=\cos\theta\,e^{-iE_+t/\hbar}|+\rangle-\sin\theta\,e^{-iE_-t/\hbar}|-\rangle.
\]

投影到 $|1\rangle$：
\[
\langle1|\psi(t)\rangle=\cos^2\theta\,e^{-iE_+t/\hbar}+\sin^2\theta\,e^{-iE_-t/\hbar}.
\]

概率：
\begin{align*}
P_{1\to1}(t)&=|\langle1|\psi(t)\rangle|^2\\
&=\cos^4\theta+\sin^4\theta+2\cos^2\theta\sin^2\theta\cos\Bigl(\frac{(E_+-E_-)t}{\hbar}\Bigr)\\
&=1-2\cos^2\theta\sin^2\theta\Bigl[1-\cos\Bigl(\frac{\Omega t}{\hbar}\Bigr)\Bigr]\\
&=1-\sin^2(2\theta)\,\sin^2\Bigl(\frac{\Omega t}{2\hbar}\Bigr).
\end{align*}

\textbf{Step 3：微扰论验证}

设 $|V|\ll\Delta=E_2-E_1$。展开精确本征值：
\[
\Omega=\Delta\sqrt{1+\frac{4|V|^2}{\Delta^2}}
\approx\Delta\Bigl(1+\frac{2|V|^2}{\Delta^2}\Bigr)=\Delta+\frac{2|V|^2}{\Delta}.
\]

故
\[
E_+=\bar{E}+\frac{\Omega}{2}\approx\bar{E}+\frac{\Delta}{2}+\frac{|V|^2}{\Delta}
=E_1+\frac{|V|^2}{E_1-E_2}.
\]

对比微扰论：$E_1^{(1)}=V_{11}$，$E_1^{(2)}=\frac{|V_{12}|^2}{E_1-E_2}$。若 $V_{11}=0$，则精确展开与微扰论二阶一致。

\textbf{Step 4：拉比振荡推导}

哈密顿量 $H=-B_z\sigma_z-\mu\sigma_x$。在 $|\uparrow\rangle,|\downarrow\rangle$ 基下：
\[
H=\begin{pmatrix}-B_z&-\mu\\-\mu&B_z\end{pmatrix}.
\]

$E_1=-B_z$, $E_2=B_z$, $V=-\mu$。故 $\Delta=2B_z$, $\bar{E}=0$。

有效分裂：
\[
\Omega=\sqrt{4B_z^2+4\mu^2}=2\sqrt{B_z^2+\mu^2}.
\]

混合角：$\tan2\theta=\frac{2\mu}{2B_z}=\frac{\mu}{B_z}$。

初态 $|\psi(0)\rangle=|\downarrow\rangle=|2\rangle$。$P_{\uparrow}(t)=P_{2\to1}(t)$：
\[
P_{\uparrow}(t)=\sin^2(2\theta)\,\sin^2\Bigl(\frac{\Omega t}{2\hbar}\Bigr)
=\frac{\mu^2}{\mu^2+B_z^2}\,\sin^2\Bigl(\frac{\sqrt{\mu^2+B_z^2}}{\hbar}t\Bigr).
\]

\textbf{Step 5：中微子振荡推导}

味本征态与质量本征态关系：
\[
\begin{pmatrix}|\nu_e\rangle\\|\nu_\mu\rangle\end{pmatrix}
=\begin{pmatrix}\cos\theta&\sin\theta\\-\sin\theta&\cos\theta\end{pmatrix}
\begin{pmatrix}|\nu_1\rangle\\|\nu_2\rangle\end{pmatrix}.
\]

$t=0$ 时 $|\nu(0)\rangle=|\nu_e\rangle=\cos\theta|\nu_1\rangle+\sin\theta|\nu_2\rangle$。

时间演化：$|\nu_k(t)\rangle=e^{-iE_kt/\hbar}|\nu_k\rangle$。

故
\[
|\nu(t)\rangle=\cos\theta\,e^{-iE_1t/\hbar}|\nu_1\rangle+\sin\theta\,e^{-iE_2t/\hbar}|\nu_2\rangle.
\]

投影回 $|\nu_e\rangle$：
\[
\langle\nu_e|\nu(t)\rangle=\cos^2\theta\,e^{-iE_1t/\hbar}+\sin^2\theta\,e^{-iE_2t/\hbar}.
\]

概率：
\begin{align*}
P_{e\to e}(t)&=|\langle\nu_e|\nu(t)\rangle|^2\\
&=1-\sin^2(2\theta)\,\sin^2\Bigl(\frac{(E_2-E_1)t}{2\hbar}\Bigr).
\end{align*}

极端相对论近似 $E_k\approx p+\frac{m_k^2}{2p}$，$E_2-E_1\approx\frac{\Delta m^2}{2E}$。

飞行距离 $L\approx ct$，故：
\[
P_{e\to e}(L)=1-\sin^2(2\theta)\,\sin^2\Bigl(\frac{\Delta m^2\,L}{4E\hbar c}\Bigr).
\]
\end{derivationbox}

\vspace{1em}
